\subsection*{What an arithmetic-location-insensitive certificate can and cannot mean}

There is no unrestricted theorem that every location-insensitive estimate
fails. The relevant information loss must be specified. Two precise
obstructions follow below: one concerns scalar density relaxations for
the actual multiplier, and one concerns its orientation relative to the
physical Fourier space. They do not identify all the earlier closures
with one hypothesis.

Put $a=\log\lambda$, $I_a=(-a,a)$, and let $\mathcal F_a$ be zero
extension followed by the unitary Fourier transform. Retain the
first-slot-linear convention and the exact $\beta_a$ in
\eqref{eq:v135-full-symbol}. Its form domain is
\[
 \mathcal D_a=\left\{f\in L^2(I_a):
 \int_{\mathbb R}\log(2+|\xi|)|\mathcal F_af(\xi)|^2d\xi<\infty\right\}.
\]
Indeed $\beta_a(\xi)=\log(2+|\xi|)+O_a(1)$; this is the same
complete physical form domain, not a new source norm. There is no
Fourier cutoff in this subsection. Either use the full space or use
the actual source complement $Q_aL^2(I_a)$, where
$Q_a=I-|u_a\rangle\langle u_a|$ and the established complete
residual $\norm{\mathsf W_{e^a}u_a}\to0$ is retained. Both parities
are included. No equality with $C_a^{\rm low}$ is asserted.

For an ordinary unit vector $f$, set
\[
 p_f(\xi)=\tfrac12\bigl(|\mathcal F_af(\xi)|^2+
                               |\mathcal F_af(-\xi)|^2\bigr).
\]
Evenness of $\beta_a$ gives $q_a[f]=\int\beta_ap_f$. The elementary
physical constraints are
\begin{equation}\label{eq:v139-density-constraints}
 p_f\ge0,\qquad \int p_f=1,\qquad
 \norm{p_f}_\infty\le\frac a\pi,\qquad
 \operatorname{TV}(p_f)\le2a.
\end{equation}
The height bound is Cauchy--Schwarz on an interval of length $2a$.
For $F=\mathcal F_af$, Plancherel gives
$\norm{F'}_2=\norm{xf}_2\le a$, and
$\norm{(|F|^2)'}_1\le2\norm{F'}_2\norm F_2\le2a$.
Reflection averaging preserves these bounds.

\begin{theorem}[Obstruction for relaxations admitting the cutoff probe]
\label{thm:v139-probe-relaxation}
Let $\mathcal R_a$ be a collection of even nonnegative probability
densities containing every $p_f$ from the full physical space, or
from its complete source complement as just specified. Assume that,
for all sufficiently large $a$, it also contains
\begin{equation}\label{eq:v139-probe-density}
 p_a^\star(\xi)=\frac a{2\pi}
 \left\{w\bigl(2a(\xi-\gamma_0)\bigr)
       +w\bigl(2a(-\xi-\gamma_0)\bigr)\right\},
\end{equation}
where $w$ is the fixed probe of Lemma~\ref{lem:v137-probe} and
$\gamma_0$ is any fixed critical-zero ordinate. If $B_a\ge0$ is a
sound relaxed error, meaning
\begin{equation}\label{eq:v139-relaxed-error}
 \int\beta_a p\ge-B_a\qquad(p\in\mathcal R_a),
\end{equation}
then $B_a\to+\infty$. In particular such a relaxation cannot supply
even a uniform finite lower floor along a cofinal family.

One explicit eligible relaxation is the set of all even $p\in
W^{1,1}(\mathbb R)$ satisfying \eqref{eq:v139-density-constraints}.
Thus retaining mass, height and total variation, with their displayed
physical constants, still loses information essential to this route.
Under RH the sharper conclusion $B_a\ge a/(6\pi^2)$ holds eventually;
that rate is not asserted unconditionally.
\end{theorem}
\begin{proof}
The unsymmetrized density $(a/\pi)w(2a(\xi-\gamma_0))$ has mass one,
height at most $a/\pi$, and variation $2a/\pi\le2a$.
Symmetrization preserves those inequalities and gives
\eqref{eq:v139-probe-density}. Compact smoothness makes its pairing
with $\beta_a$ finite. By evenness and
Proposition~\ref{prop:v137-conditional-mass}, under RH
\[
 \int\beta_ap_a^\star=\frac a\pi J_a
       \le-\frac a{6\pi^2}
\]
for all sufficiently large $a$. This proves the conditional rate.
If $B_{a_j}$ were bounded on any cofinal sequence, inclusion of all
physical densities would give a complete uniform lower floor. In
the source-complement case apply
Corollary~\ref{cor:v136-complement-floor}; in the full-space case
apply Proposition~\ref{prop:v136-bounded-floor}. Either implication
gives RH, and hence contradicts the just proved conditional rate.
The absence of a bounded cofinal subsequence proves the asserted
unconditional divergence.
\end{proof}

The quantifier over both parities is essential to this proof.
For separate parity errors it proves unconditional divergence of
their maximum, not unconditional divergence of each merely from
one sector's boundedness. The probe is not asserted to be a physical
squared transform or to satisfy the sharper source-dependent
pointwise evaluation bounds. In fact a nonzero compactly supported
frequency density cannot be such a squared transform: the entire
transform of a compact physical function cannot vanish on a real
open interval. The obstruction is precisely an excess of tests
in the relaxation, not a negative physical Weil test.

There is a particularly explicit location-insensitivity property.
For a real symbol $b$ bounded below with finite-measure sublevel sets,
define $b^\uparrow(s)$ to be its increasing rearrangement with respect
to Lebesgue measure. The strongest universal lower bound using only
unit mass and the height bound $a/\pi$ is
\begin{equation}\label{eq:v139-bathtub}
 \ell_a^{\rm cap}(b):=
 \inf_{\substack{p\ge0,\ \int p=1\\p\le a/\pi}}
        \int bp
 =\frac a\pi\int_0^{\pi/a} b^\uparrow(s)\,ds.
\end{equation}
This is unchanged by any measure-preserving rearrangement of $b$:
it retains the full distribution of favorable and unfavorable values
but discards their physical frequency positions. The formula is the
classical bathtub principle~\cite[Section 1.14]{LiebLoss2001}.
For clarity, choose a threshold $t$ whose lower sublevel has measure
at most $\pi/a$ and whose closed sublevel has measure at least
$\pi/a$, and fill that lowest region with density $a/\pi$, partially
filling the level $t$ if needed. For every competitor,
$(b-t)(p-p_*)\ge0$ pointwise; its integral and equal total masses
prove minimality and \eqref{eq:v139-bathtub}.
For even $b$, symmetrizing does not change the value.
Theorem~\ref{thm:v139-probe-relaxation} implies
\begin{equation}\label{eq:v139-cap-divergence}
 \ell_a^{\rm cap}(\beta_a)\longrightarrow-\infty.
\end{equation}
Any lower-bound rule valid for \emph{every} mass-and-height-admissible
density is at most \eqref{eq:v139-bathtub}. This is the exact scope of
the distribution-only assertion; it does not prohibit every possible
theorem whose input happens to include a histogram.

The scalar primitive criterion is included by a different verified
hypothesis: its decomposition $b=c+V'$, $c\ge0$,
$\norm V_\infty=\Delta(b)/2$, is valid against every nonnegative
unit density of variation at most $2a$. Thus its error $a\Delta(b)$
satisfies \eqref{eq:v139-relaxed-error}. Notice that $\Delta(b)$
itself is not rearrangement invariant. Calling it distribution-only
would be incorrect.

\begin{theorem}[Spectral-orientation obstruction with protected finite data]
\label{thm:v139-protected-orbit}
Fix $a$, let $M$ be multiplication by the exact $\beta_a$ on
$L^2(\mathbb R)$ with its multiplication-operator domain, and let
$S$ be any finite-dimensional frequency subspace. Let $\mathcal U_S$
consist of unitaries which fix $S$ pointwise, differ from the identity
by finite rank, and preserve $\mathcal D(M)$. Then
\begin{equation}\label{eq:v139-orbit-edge}
 \inf_{U\in\mathcal U_S}
 \inf_{\substack{f\in C_c^\infty(I_a),\ \norm f=1\\
                 \mathcal F_af\perp S}}
 \ip{U^*MU\mathcal F_af}{\mathcal F_af}
       =\operatorname*{ess\,inf}\beta_a.
\end{equation}
For any finite physical subspace $E$ with
$\mathcal F_aE\subset\mathcal D(M)$, one can take
$S\supset\mathcal F_aE+M\mathcal F_aE$. Every allowed unitary then
preserves the complete operator action on the protected columns:
\begin{equation}\label{eq:v139-protected-columns}
 U^*MU\mathcal F_ae=M\mathcal F_ae\qquad(e\in E).
\end{equation}
Consequently a sound universal compression bound insensitive to all
the allowed orientations, even after retaining these exact columns,
cannot exceed $\operatorname*{ess\,inf}\beta_a$. Along the actual
arithmetic family that upper limit tends to $-\infty$.
The result holds in either parity, with parity-preserving unitaries
and reflection-invariant protected data.
\end{theorem}
\begin{proof}
Write $m=\operatorname*{ess\,inf}\beta_a$. The lower bound in
\eqref{eq:v139-orbit-edge} is immediate from $M\succeq mI$.
Fix $r>m$. The spectral subspace $1_{\{\beta_a<r\}}L^2$ is
infinite dimensional because its support has positive Lebesgue measure;
its vectors lie in $\mathcal D(M)$ since $m\le\beta_a<r$ there.
Choose a unit vector $v$ in this subspace orthogonal to $S$.
Choose a unit smooth physical $f$ satisfying the finitely many
conditions $\mathcal F_af\perp S$, and put $h=\mathcal F_af$.
Such $f$ exists because the smooth physical space is infinite
dimensional. Its rapidly decreasing transform belongs to
$\mathcal D(M)$. On the space $\operatorname{span}\{h,v\}\subset
S^\perp$, choose a unitary carrying $h$ to $v$, and extend it by the
identity. Then $U-I$ has finite rank with range in $\mathcal D(M)$;
both $U$ and $U^*$ preserve that domain. Its Rayleigh value on $h$
is $\ip{Mv}v<r$. Let $r\downarrow m$.

If $S$ contains both protected columns and their images, $U$ fixes
each of them, which proves \eqref{eq:v139-protected-columns}.
For the constructed unitary, expansion of $U=I+K$ shows that
$U^*MU-M=K^*M+MK+K^*MK$ extends to a bounded finite-rank
self-adjoint perturbation: $MK$ is bounded and $K^*M$ is its
adjoint on the original domain. Thus the compressed forms also
keep their original domain. Evenness of $\beta_a$ leaves infinite
dimensional negative spectral subspaces in each parity, so the same
construction inside either sector proves the parity statement.
Finally apply Corollary~\ref{cor:v137-scalar-no-go}.
\end{proof}

For the actual finite Fourier source proxy, the additional domain
condition is automatic: its zero extension is piecewise smooth,
its transform is $O_a((1+|\xi|)^{-1})$, and multiplication by
$\beta_a=O_a(\log(2+|\xi|))$ remains square integrable. Protecting
it therefore preserves its exact compressed source residual, not
just the source norm or its Rayleigh value. Any finite number of
other domain columns can also be protected, even if their number
grows with $a$. This does not preserve the action on an infinite
physical block.

Spectral compression and orientation are classical concerns; see
Fan--Pall~\cite{FanPall1957}. The project-specific extension here is
the explicitly protected graph data and its application to the
actual divergent symbol. Most importantly, $U^*MU$ generally is
\emph{not} a multiplication operator. The theorem rules out bounds
required to be valid on that larger operator class. It does not
construct a rearranged prime sequence, a negative arithmetic Weil
test, or a no-go for methods retaining the multiplier's physical
support geometry. It is not a new finite-rank metric repair.

\begin{proposition}[Why mere location independence is too broad]
\label{prop:v139-location-counterexample}
Let $D_a>0$, choose any measurable set $E_a$ of measure
$\pi/(2a(D_a+1))$, and put
$b_a=1-(D_a+1)1_{E_a}$. Then for every physical $f$,
\[
 \int b_a|\mathcal F_af|^2\ge\tfrac12\norm f^2,
\]
independently of the location or shape of $E_a$, even if
$D_a\to\infty$ and hence $\operatorname*{ess\,inf}b_a\to-\infty$.
\end{proposition}
\begin{proof}
Use $|\mathcal F_af|^2\le(a/\pi)\norm f^2$ and integrate over
$E_a$. The subtracted contribution is at most
$(D_a+1)(a/\pi)|E_a|\norm f^2=\norm f^2/2$.
\end{proof}

This example is not the arithmetic symbol. It proves that neither
negative depth alone nor the informal phrase ``ignores arithmetic
locations'' implies the requested universal conclusion. The actual
symbol obstruction in Theorem~\ref{thm:v139-probe-relaxation} uses
the signed cutoff-probe pairing, and not just its minimum.
The proved scalar pointwise and primitive closures fall within the
displayed classes. The Schatten, signed block-norm, noncommuting
band, Cotlar, and flat-top results have additional geometric or
structural hypotheses; their recorded closures are not corollaries
of a single insensitivity assertion. Their precise coverage is
tabulated in the new evidence report. The surviving joint signed
concentration operator retains information excluded by both
relaxations. Its cofinal finite-floor bound is still unproved.
All endpoint, prime-cut, logarithmic-diagonal, Fourier normalization
and sampler graph-defect conventions remain unchanged. No G2 sign
gap is closed, and no worldwide novelty or publication readiness is
asserted.
